\documentclass[11pt]{article}
\usepackage[a4paper,margin=22mm]{geometry}
\usepackage[no-math]{fontspec}
\setmainfont{lmroman10-regular.otf}[BoldFont=lmroman10-bold.otf,ItalicFont=lmroman10-italic.otf,BoldItalicFont=lmroman10-bolditalic.otf]
\usepackage{amsmath,amssymb,amsthm,mathtools,graphicx,xcolor,hyperref,xurl,xspace,ifthen}
\setlength{\emergencystretch}{4em}
\SetMathAlphabet{\mathbf}{normal}{OT1}{cmr}{bx}{n}
\SetMathAlphabet{\mathbf}{bold}{OT1}{cmr}{bx}{n}
\newtheorem{theo}{Theorem}[section]
\newtheorem{lemm}[theo]{Lemma}
\newtheorem{lem}[theo]{Lemma}
\newtheorem{lemma}[theo]{Lemma}
\newtheorem{prop}[theo]{Proposition}
\newtheorem{coro}[theo]{Corollary}
\newtheorem{corol}[theo]{Corollary}
\newtheorem{conj}[theo]{Conjecture}
\newtheorem{Proposition}[theo]{Proposition}
\newtheorem{theoreme}[theo]{Theorem}
\theoremstyle{definition}
\newtheorem{defi}[theo]{Definition}
\newtheorem{exam}[theo]{Example}
\newtheorem{exem}[theo]{Example}
\newtheorem{rema}[theo]{Remark}
\newtheorem{nota}[theo]{Notation}
\newtheorem*{rema*}{Remark}
\newtheorem*{defi*}{Definition}
\newtheorem*{theo*}{Theorem}
\newtheorem*{lemm*}{Lemma}
\newtheorem*{prop*}{Proposition}
\newcommand{\enoncename}{}
\newtheorem{corpusenonce}[theo]{\enoncename}
\newenvironment{enonce}[1]{\renewcommand{\enoncename}{#1}\begin{corpusenonce}}{\end{corpusenonce}}
\newenvironment{enonce*}[1]{\par\medskip\noindent\textbf{#1.}\itshape}{\par\medskip}
\newenvironment{proofc}{\begin{proof}}{\end{proof}}
\newcommand{\Mc}{\mathcal{M}}
\newcommand{\pointir}{.\ ---\ }
\providecommand{\mathof}[1]{\mathrm{#1}}

\numberwithin{equation}{section}
\numberwithin{figure}{section}
\numberwithin{table}{section}
\def\endappendix{}
%~Mouliné par MaN_auto v.0.23.0 2020-08-13 13:19:33


\newfontfamily\CorpusDoubleStroke{Noto Sans Math}
\newcommand{\mathds}[1]{{\ifthenelse{\equal{#1}{1}}{\mathord{\text{\CorpusDoubleStroke\char"1D7D9}}}{\mathbb{#1}}}}
%\usepackage{ifthen}

\DeclareMathOperator{\im}{Im}
\DeclareMathOperator{\ind}{\mathds{1}}





%\AtBeginDocument{ }



\newcommand{\pr}{\mathds{P}}
\newcommand{\CC}{\mathds{C}}
\newcommand{\R}{\mathds{R}}
\newcommand{\ph}{\varphi}
\newcommand{\bb}{{\mathrm{b}}}
\newcommand{\oo}{{\mathrm{o}}}
\newcommand{\si}{{\mathrm{s}}}

%\renewcommand{\le}{\leqslant}
\renewcommand{\ge}{\geqslant}


\newcommand{\formula}[2][nolabel] {\ifthenelse{\equal{#1}{nolabel}} {\begin{align*}
#2
\end{align*}}
{\ifthenelse{\equal{#1}{}} {\begin{align}
#2
\end{align}}
{\begin{align}\label{#1}
#2
\end{align}}}}


%%%%%%%%%%%%%%%%%%%%%%%%%%%%%%%%%%%%%%%%%%
\graphicspath{{./figures/}}

\newcommand*{\mk}{\mkern -1mu}
\newcommand*{\Mk}{\mkern -2mu}
\newcommand*{\mK}{\mkern 1mu}
\newcommand*{\MK}{\mkern 2mu}

\hypersetup{urlcolor=purple, linkcolor=blue, citecolor=red}

\let\oldhat\hat
\renewcommand*{\hat}[1]{\mathchoice{\widehat{#1}}{\widehat{#1}}{\oldhat{#1}}{\oldhat{#1}}}
%%%%%%%%%%%%%%%%%%%%%%%%%%%%%%%%%%%%%%%%%

\begin{document}
\section*{\raggedright 2642: Determination of nondegenerate finite signed measures by all convolution restrictions to the positive half-line}
\noindent\textbf{Author:} Mateusz Kwaśnicki\par\noindent\textbf{Source:} Random walks are determined by their trace on the positive half-line. Annales Henri Lebesgue 3 (2020), publisher version, DOI 10.5802/ahl.64. \url{https://ahl.centre-mersenne.org/item/AHL_2020__3__1389_0/}
\par\noindent\textbf{License:} CC BY 4.0, \url{https://creativecommons.org/licenses/by/4.0/}. The original publisher PDF has the explicit license notice. See the saved source/license records.
\par\noindent\textbf{Editorial material note (not author text):} The selected problem is the finite signed-measure assertion of Theorem 1.1: nonzero restrictions to the positive half-line and equality there of every positive-integer convolution power determine the entire measures. The complete author body retains the local induction, holomorphic gluing, simultaneous all-powers factorization and removability argument. Random-walk, ladder and Levy reformulations and the final counterexample discussion are unselected context. All original author labels and formulas remain unchanged. All proof text and formulas below are editable author TeX. Mathematical wording, language and proof order are retained. Only the article class, title matter and bibliography presentation are adapted for independent compilation; no proof has been generated or repaired.
\par\noindent\textbf{Editorial difficulty:} H1: The proof turns convolution identities into simultaneous bounded-type factorizations, controls singular measures through every power, and removes boundary singularities. This resolves a research uniqueness problem beyond ordinary qualifying analysis.
\par\noindent\textbf{Editorial retrieval and attribution:} Retrieved 2026-10-01. DOI 10.5802/ahl.64. Source attribution notice: ../licenses/2642-source-license.txt; complete legal text: ../licenses/CC-BY-4.0.txt. Original source excerpts and the existing source-material check are retained in ../sources/2642/.
\medskip\hrule\medskip
\noindent\textbf{Author statement, complete proof and local supporting text follow in source order.}
\section{Introduction and main result}

In this note we give an affirmative answer to a question posed by Lo\"ic Chaumont and Ron Doney in~\cite{bib:cd17}, inspired by Vincent Vigon's conjecture in~\cite{bib:v02}. The main result was previously stated without proof in a more general form in~\cite{bib:ou90}, and an erroneous proof was given in~\cite{bib:u94}.

A random walk $X_n$ is said to be \emph{non-degenerate} if $\pr(X_n > 0) \ne 0$. Similarly, a finite signed Borel measure $\mu$ on $\R$ is said to be \emph{non-degenerate} if the restriction of $\mu$ to $(0, \infty)$ is a non-zero measure.

\begin{theo}\label{thm:wh}
If $X_n$ and $Y_n$ are non-degenerate random walks such that $\max(X_n, 0)$ and $\max(Y_n, 0)$ are equal in distribution for all $n = 1, 2,\,\ldots\,$, then $X_n$ and $Y_n$ are equal in distribution for $n = 1, 2,\,\ldots$

More generally, if $\mu$ and $\nu$ are non-degenerate finite signed Borel measures and their $n$-fold convolutions $\mu^{*n}$ and $\nu^{*n}$ agree on $(0, \infty)$ for $n = 1, 2,\,\ldots\,$, then $\mu = \nu$.
\end{theo}

Following~\cite{bib:cd17}, we remark that various reformulations of the above result are possible. A non-degenerate random walk $X_n$ is determined by any of the following objects:
\begin{itemize}
\item The law of the ascending ladder process $(T_k, S_k)$; here $S_k = X_{T_k}$ is the $k^{\rm th}$ running maximum of the random walk.
\item The upward space-time Wiener--Hopf factor $\Phi_+(q, \xi)$, that is, the characteristic function of $(T_1, S_1)$.
\item The distributions of the running maxima $\max(0, X_1, X_2,\,\ldots,\,X_n)$ for all\linebreak$n = 1, 2,\,\ldots$
\end{itemize}
Theorem~\ref{thm:wh} clearly implies that a~non-degenerate Lévy process $X_t$ is determined by any of the following objects:
\begin{itemize}
\item The distributions of $\max(X_t, 0)$ for all $t > 0$ (or even for $t = 1, 2,\,\ldots$).
\item The law of the ascending ladder process $(T_t, S_t)$.
\item The upward space-time Wiener--Hopf factor $\kappa_+(q, \xi)$, that is, the characteristic exponent of $(T_t, S_t)$.
\item The distributions of the running suprema $\sup\{X_s : s \in [0, t]\}$ for all $t > 0$.
\end{itemize}
For further discussion, we again refer to~\cite{bib:cd17}, where Theorem~\ref{thm:wh} was proved under various relatively mild additional conditions. For related research, see~\cite{bib:cd17,bib:lms76,bib:o85,bib:ou90,bib:u90,bib:u94} and the references therein.

Theorem~\ref{thm:wh} was given without proof in~\cite{bib:ou90} in a more general form: Theorem~4 therein claims that $\mu = \nu$ if $\mu$ and $\nu$ are non-degenerate finite Borel measures on $\R$ and the restrictions of $\mu^{*n_k}$ and $\nu^{*n_k}$ to $(0, \infty)$ are equal for $k = 1, 2,\,\ldots\,$, where $n_1 = 1$ and $n_2 - 1, n_3 - 1,\,\ldots$ are distinct and have no common divisor other than~$1$. Noteworthy, this result is stated for measures on the Euclidean space of arbitrary dimension, and their restrictions to the half-space. A proof is given in~\cite{bib:u94} under the additional condition $n_2 = 2$, and only in dimension one. However, the argument in~\cite{bib:u94} contains a gap, that we describe at the end of this article.


\section{Proof}

All measures considered below are finite, signed Borel measures. For a measure $\mu$ on $\R$, we denote the restrictions of $\mu$ to $(0, \infty)$ and $(-\infty, 0]$ by $\mu_+ = \ind_{(0, \infty)} \mu$ and $\mu_- = \ind_{(-\infty, 0]} \mu$. This should not be confused with the Hahn decomposition of $\mu$ into the positive and negative part. By $\mu^{*n}$ we denote the $n$-fold convolution of $\mu$, and we define $\mu^{*0}$ to be the Dirac measure $\delta_0$. For brevity, we write $\mu_\pm^{*n} = (\mu_\pm)^{*n}$, as opposed to $(\mu^{*n})_\pm$. We record the following elementary identities: $(\delta_0 * \sigma_-)_+ = (\sigma_-)_+ = 0$, $(\pi_- * \sigma_-)_+ = 0$, and $(\pi * \sigma_-)_+ = (\pi_+ * \sigma_-)_+$.

We denote the characteristic function of a measure $\mu$ by $\hat\mu$:
\[
\hat\mu(z) = \int_\R e^{i z x} \mu(dx) 
\]
for $z \in \R$, and also for those $z \in \CC$ for which the integral converges. We recall that $\hat\mu_+$ is a bounded holomorphic function in the upper complex half-plane $\CC_+ =\linebreak\{z \in \CC : \im z > 0\}$, continuous on the boundary. Similarly, $\hat\mu_-$ is a bounded holomorphic function on the lower complex half-plane $\CC_- = \{z \in \CC : \im z < 0\}$.

\begin{lemm}\label{lem:mupm}
Suppose that $\mu, \nu$ are measures on $\R$ satisfying
\[
(\mu^{*n})_+ = (\nu^{*n})_+\quad\text{for }n = 1, 2,\,\ldots,\,N.
\]
Then $\mu_+ = \nu_+$ and
\begin{equation}\label{eq:mupm}
\left(\mu_+^{*n} * \mu_-^{*k}\right)_+ = \left(\nu_+^{*n} * \nu_-^{*k}\right)_+ \quad\text{for }n = 0, 1,\,\ldots,\,N - 1\text{ and }k = 1, 2,\,\ldots
\end{equation}
\end{lemm}

\begin{proof}
We proceed by induction with respect to $N$. For $N = 1$ the result is trivial: we have $\mu_+ = (\mu^{*1})_+ = (\nu^{*1})_+ = \nu_+$ and $(\mu_+^{*0} * \mu_-^{*k})_+ = (\delta_0 * \mu_-^{*k})_+ = 0\linebreak= (\delta_0 * \nu_-^{*k})_+ = (\nu_+^{*0} * \nu_-^{*k})_+$ for $k = 1, 2,\,\ldots\,$ Suppose that the assertion of Lemma~\ref{lem:mupm} holds for some $N$, and suppose that $(\mu^{*n})_+ = (\nu^{*n})_+$ for $n = 1, 2,\,\ldots,\,N,\linebreak N + 1$. By the induction hypothesis, formula~\eqref{eq:mupm} holds for $n = 0, 1,\,\ldots,\,N - 1$ and $k = 1, 2, \ldots\,$, and we have $\mu_+ = \nu_+$. Therefore, we only need to prove~\eqref{eq:mupm} for $n = N$ and $k = 1, 2, \ldots$

By the binomial theorem,
\begin{align*}
0 & = \left(\mu^{*N + 1} - \nu^{*N + 1}\right)_+ = \left((\mu_+ + \mu_-)^{*N + 1} - (\nu_+ + \nu_-)^{*N + 1}\right)_+ \\
& = \sum_{j = 0}^{N + 1} \binom{N + 1}{j} \left(\mu_+^{*j} * \mu_-^{*N + 1 - j} - \nu_+^{*j} * \nu_-^{*N + 1 - j}\right)_+. 
\end{align*} 
We already know that $\mu_+^{*N + 1} = \nu_+^{*N + 1}$ and $(\mu_+^{*j} * \mu_-^{*N + 1 - j})_+ = (\nu_+^{*j} * \nu_-^{*N + 1 - j})_+$ for $j = 1, 2,\,\ldots,\,N - 1$. Furthermore, $(\mu_-^{*N + 1})_+ = 0 = (\nu_-^{*N + 1})_+$. It follows that all terms corresponding to $j \ne N$ in the above sum are zero. Thus,
\[
0 = \binom{N + 1}{N} \left(\mu_+^{*N} * \mu_- - \nu_+^{*N} * \nu_-\right)_+,
\]
which proves~\eqref{eq:mupm} for $n = N$ and $k = 1$. The proof for $n = N$ and $k > 1$ proceeds again by induction. Suppose that~\eqref{eq:mupm} holds for $n = N$ and $k = 1, 2, \ldots, K$. By the identity $(\pi * \sigma_-)_+ = (\pi_+ * \sigma_-)_+$,
\[
\left(\mu_+^{*N} * \mu_-^{*K + 1}\right)_+ = \left(\mu_+^{*N} * \mu_-^{*K} * \mu_-\right)_+ = \left(\left(\mu_+^{*N} * \mu_-^{*K}\right)_+ * \mu_-\right)_+.
\]
Applying~\eqref{eq:mupm}, with $n = N$ and $k = K$, and then again the identity $(\pi * \sigma_-)_+\linebreak= (\pi_+ * \sigma_-)_+$, we find that
\[
\left(\left(\mu_+^{*N} * \mu_-^{*K}\right)_+ * \mu_-\right)_+ = \left(\left(\nu_+^{*N} * \nu_-^{*K}\right)_+ * \mu_-\right)_+ = \left(\nu_+^{*N} * \nu_-^{*K} * \mu_-\right)_+.
\]
Recall that $\mu_+ = \nu_+$, so that
\[
\left(\nu_+^{*N} * \nu_-^{*K} * \mu_-\right)_+ = \left(\mu_+^{*N} * \mu_- * \nu_-^{*K}\right)_+.
\]
We use the identity $(\pi * \sigma_-)_+ = (\pi_+ * \sigma_-)_+$ for the third time, and then we again apply~\eqref{eq:mupm}, with $n = N$ and $k = 1$:
\[
\left(\mu_+^{*N} * \mu_- * \nu_-^{*K}\right)_+ = \left(\left(\mu_+^{*N} * \mu_-\right)_+ * \nu_-^{*K}\right)_+ = \left(\left(\nu_+^{*N} * \nu_-\right)_+ * \nu_-^{*K}\right)_+.
\]
Finally, once again we apply the identity $(\pi * \sigma_-)_+ = (\pi_+ * \sigma_-)_+$:
\[
\left(\left(\nu_+^{*N} * \nu_-\right)_+ * \nu_-^{*K}\right)_+ = \left(\nu_+^{*N} * \nu_- * \nu_-^{*K}\right)_+ = \left(\nu_+^{*N} * \nu_-^{*K + 1}\right)_+.
\]
The above chain of equalities implies that $(\mu_+^{*N} * \mu_-^{*K + 1})_+ = (\nu_+^{*N} * \nu_-^{*K + 1})_+$, which is just~\eqref{eq:mupm} with $n = N$ and $k = K + 1$. We conclude that~\eqref{eq:mupm} holds for $n = N$ and every $k = 1, 2,\,\ldots\,$, and the proof of Lemma~\ref{lem:mupm} is complete.
\end{proof}

A holomorphic function $f$ on $\CC_-$ is said to be of \emph{bounded type} (or belong to the \emph{Nevanlinna class}) if $\log |f(x)|$ has a harmonic majorant on $\CC_-$. Equivalently, $f$ is of bounded type if it is a ratio of two bounded holomorphic functions on $\CC_-$. We recall the following fundamental factorisation theorem for holomorphic functions on $\CC_-$ which are bounded or of bounded type, and we refer to~\cite{bib:g07,bib:m09} for further details.

\begin{theo}{\cite[Theorem~II.5.5 and Corollary~II.5.7]{bib:g07}; \cite[Theorem~13.15]{bib:m09}}\label{thm:fact}

Let $f$ be a holomorphic function of bounded type on the lower complex half-plane, and suppose that $f$ is not identically zero. Let $\alpha_0$ be the multiplicity of the zero of $f$ at $z = -i$ (possibly $\alpha_0 = 0$), and let $z_1, z_2, \ldots$ be the (finite or infinite) sequence of all zeros of $f$ in the lower complex half-plane, with corresponding multiplicities $\alpha_1, \alpha_2,\,\ldots\,$ Then $f$ admits a factorisation
\begin{equation}\label{eq:fact}
f(z) = f_\bb(z) f_\oo(z) f_\si(z)
\end{equation}
(unique, up to multiplication of $f_\oo$ and $f_\si$ by a constant of modulus~$1$), with the following factors. The function $f_\bb$ is a \emph{Blaschke product}, determined uniquely by the zeros of $f$:
\begin{equation}\label{eq:bb}
f_\bb(z) = \biggl(\frac{z + i}{z - i} \biggr)^{\!\alpha_0} \prod_j \biggl(\frac{|1 + z_j^2|}{1 + z_j^2} \, \frac{z - z_j}{z - \bar{z}_j} \biggr)^{\!\alpha_j}.
\end{equation}
The function $f_\oo$ is an \emph{outer function}, a holomorphic function determined uniquely up to multiplication by a constant of modulus~$1$ by the formula:
\begin{equation}\label{eq:oo}
|f_\oo(z)| = \exp \biggl(\frac{1}{\pi} \int_{-\infty}^\infty \frac{-\im z}{|z - x|^2} \, \log |f(x)| \, dx \biggr).
\end{equation}
Finally, the function $f_\si$ is a \emph{singular inner function}, a holomorphic function determined uniquely up to multiplication by a constant of modulus~$1$ by the expression:
\begin{equation}\label{eq:si}
|f_\si(z)| = \exp \biggl(a \im z - \frac{1}{\pi} \int_\R \frac{-\im z}{|z - x|^2} \, \sigma(dx) \biggr),
\end{equation}
where $a \in \R$ is a constant and $\sigma$ is a signed measure, singular with respect to the Lebesgue measure.

Furthermore, for almost all $x \in \R$ with respect to both the Lebesgue measure and the measure $\sigma$, the limit $f(x)$ of $f(x + i y)$ as $y \to 0^-$ exists. This boundary limit $f(x)$ is non-zero almost everywhere with respect to the Lebesgue measure and zero almost everywhere with respect to $\sigma$. The symbol $f(x)$ used in the definition of the outer function $f_\oo$ refers precisely to this boundary limit. Additionally, we have \[
\sum_j \alpha_j |\im z_j| \left(1 + |z_j|^2\right)^{-1} < \infty, \int_{-\infty}^\infty \left(1 + x^2\right)^{-1} |\log |f(x)|| dx < \infty
\] 
and
\[
\int_\R \left(1 + x^2\right)^{-1} |\sigma|(dx) < \infty\,,
\]
 and any parameters $\alpha_j$, $z_j$, $a$, $\sigma$ and boundary values $|f(x)|$, $x \in \R$, which satisfy these conditions, correspond to some function $f$ of bounded type.

Finally, $f$ is a bounded holomorphic function in the lower complex half-plane if and only if $a \ge 0$, $\sigma$ is a non-negative measure and the boundary values $|f(x)|$ are bounded for $x \in \R$.
\end{theo}

\begin{lemm}\label{lem:ext}
Suppose that $\mu$ is a measure on $\R$ such that $\mu_-$ is a non-zero measure and $(\mu_+ * \mu_-)_+ = 0$. Then $\hat\mu_+$ has a holomorphic extension $\ph$ to the connected open set
\[
D = \CC \setminus \left\{z \in \CC_- \cup \R : \hat\mu_-(z) = 0\right\},
\]
and $\ph$ is a meromorphic function on $\CC \setminus \{z \in \R : \hat\mu_-(z) = 0\}$. Furthermore, $\ph \hat\mu_-$ extends to a function which is holomorphic on $\CC_-$ and continuous on $\CC_- \cup \R$, namely, the characteristic function of $\mu_+ * \mu_-$.
\end{lemm}

\begin{proof}
Denote $\nu = \mu_+ * \mu_-$; by the assumption, $\nu = \nu_-$. Let $f = \hat\mu_+$, $g = \hat\mu_-$ and $h = \hat \nu = \hat \nu_-$. Clearly, $h(z) = f(z) g(z)$ for $z \in \R$. Let
\[
A = \{ z \in \R : g(z) = 0 \},\quad B = \{ z \in \CC_- : g(z) = 0 \},
\]
so that $D = \CC \setminus (A \cup B)$.

We note basic properties of $A$ and $B$. By continuity of $g$, $A$ and $A\cup B$ are closed sets, and $D$ is an open set. Since $g$ is holomorphic on $\CC_-$ (and not identically zero), $B$ is a countable (possibly finite) set with no accumulation points on $\CC_-$. By Theorem~\ref{thm:fact}, $A$ has zero Lebesgue measure (as a subset of $\R$). In particular, $D$ is connected. Indeed: the sets $D \cap \CC_+ = \CC_+$ and $D \cap \CC_- = \CC_- \setminus B$ are clearly path-connected, the set $D \cap \R = \R \setminus A$ is non-empty, and since $D$ is open, each point of $D \cap \R$ is path-connected with points from both $D \cap \CC_+$ and $D \cap \CC_-$.

We define a function $\ph$ on $D$ by the formula
\[
\ph(z) =
\begin{cases}
f(z) & \text{if }z \in \CC_+ \cup (\R \setminus A), \\[0.5em]
\dfrac{h(z)}{g(z)} & \text{if }z \in \CC_- \setminus B\,.
\end{cases}
\]
By definition, $\ph$ is holomorphic both on $\CC_+$ and on $\CC_- \setminus B$, as well as meromorphic on $\CC_-$. Furthermore, $\ph$ is continuous at each point $z \in \R \setminus A$, because both $f$ (defined on $\CC_+ \cup \R$) and $h / g$ (defined on $(\CC_- \setminus B) \cup (\R \setminus A)$) are continuous at\linebreak$z$ and $f(z) = h(z) / g(z)$. By a standard application of Morera's theorem (see~\cite[Theorem~IV.5.10 and Exercise~IV.5.9]{bib:c73}, or~\cite[Exercise~II.12]{bib:g07}), $\ph$ is holomorphic in $D$. It remains to note that $\ph(z) g(z) = h(z)$ for $z \in \CC_- \setminus B$.
\end{proof}

\begin{lemm}\label{lem:wh}
If $\mu$ is a measure on $\R$ such that $(\mu_+^{*n} * \mu_-)_+ = 0$ for all\linebreak$n = 1, 2,\,\ldots\,$, then either $\mu_+$ or $\mu_-$ is a zero measure.
\end{lemm}

\begin{proof}
Let $\mu$ be such a measure, and suppose that both $\mu_+$ and $\mu_-$ are non-zero measures. Let $\ph, f, g, h, A, B, D$ be as in the proof of Lemma~\ref{lem:ext}. Clearly, $\ph^n$ is the holomorphic extension of $f^n$, the characteristic function of $\mu_+^{*n}$. An application of Lemma~\ref{lem:ext} to the measure $\mu_+^{*n} + \mu_-$ implies that for all $n = 1, 2, \ldots\,$, the function $\ph^n g$ extends from $\CC_- \setminus B$ to a function $h_n$ which is bounded and holomorphic on $\CC_-$ and continuous on $\CC_- \cup \R$, namely, $h_n$ is the characteristic function of $\mu_+^{*n} * \mu_-$.

Consider the factorisations $g = g_\bb\,g_\oo\,g_\si$ and $h_n = h_{n,\bb}\,h_{n,\oo}\, h_{n,\si}$ given in Theorem~\ref{thm:fact}, and let $\sigma_g$, $a_g$ and $\sigma_{h,n}$, $a_{h,n}$ denote the corresponding non-negative measures $\sigma$ and constants $a$ for $g$ and $h_n$, respectively. Note that Theorem~\ref{thm:fact} applies both to $g$ and to $h_n = \ph^n g$, as these functions are not identically zero: $f$ and $g$ are characteristic functions of non-zero measures $\mu_+$ and $\mu_-$, while $h_n$ is the product of $g$ and the holomorphic extension of $f^n$.

Recall that $\ph^n = h_n / g$ on $\CC_- \setminus B$. It follows that if $\ph_{n,\bb} = h_{n,\bb} / g_\bb$, $\ph_{n,\oo} = h_{n,\oo} / g_\oo$ and $\ph_{n,\si} = h_{n,\si} / g$, then
\[
\ph^n = \ph_{n,\bb}\,\ph_{n,\oo}\,\ph_{n,\si} \]
on $\CC_- \setminus B$. Let us examine the above factors in more detail.

By definition, $\ph_{n,\oo}$ and $\ph_{n,\si}$ have no zeros in $\CC_-$. This means that if $z_0 \in \CC_-$ is a pole of $\ph$ of order $\alpha_0$, then $z_0$ is a pole of $\ph_{n,\bb} = h_{n,\bb} / g_\bb$ of order $n \alpha_0$, and therefore $g_\bb$ has a zero at $z_0$ of multiplicity at least $n \alpha_0$ for all $n = 1, 2, \ldots$\, Since all zeroes of $g_\bb$ have finite multiplicity, $\ph$ has no poles in $\CC_-$. In particular, $\ph$ extends to a~holomorphic function on $\CC \setminus A$, which will be denoted again by $\ph$, and $\ph_{n,\bb} = h_{n,\bb} / g_\bb$ has no poles in $\CC_-$. Therefore, the zeros of $h_{n,\bb}$ must cancel the zeros of $g_\bb$, and $\ph_{n,\bb}$ is a~Blaschke product.

Since $h_n(x) / g(x) = (f(x))^n$ for $x \in \R \setminus A$ and $A$ has Lebesgue measure zero, we have
\begin{align*}
|\ph_{n,\oo}(z)| & = \exp \biggl(\frac{1}{\pi} \int_{-\infty}^\infty \frac{-\im z}{|z - x|^2} \, (\log |h_n(x)| - \log |g(x)|) \, dx \biggr) \\
& = \exp \biggl(\frac{1}{\pi} \int_{-\infty}^\infty \frac{-\im z}{|z - x|^2} \, \log |f(x)|^n \, dx \biggr). 
\end{align*}
In particular, $\ph_{n,\oo}$ is a bounded outer function, namely, the outer function in the factorisation of the bounded holomorphic function $(\overline{f(\bar{z})})^n$ on the lower complex half-plane.

Finally $\ph_{n,\si}$ is the ratio of two singular inner functions, and hence a singular inner function. If we denote $a_{\ph,n} = a_{h,n} - a_g$ and $\sigma_{\ph,n} = \sigma_{h,n} - \sigma_g$, then
\[
|\ph_{n,\si}(z)| = \exp \biggl(-a_{\ph,n} \im z - \frac{1}{\pi} \int_\R \frac{-\im z}{|z - x|^2} \, \sigma_{\ph,n}(dx) \biggr).
\]

The above properties imply that $\ph^n$ is of bounded type, and therefore the factors $\ph_{n,\bb}$, $\ph_{n,\oo}$, $\ph_{n,\si}$, the signed measure $\sigma_{\ph,n}$ and the constant $a_{\ph,n} \in \R$ are uniquely determined (up to multiplication by a constant of modulus $1$ in case of $\ph_{n,\oo}$ and~$\ph_{n,\si}$).

By comparing the factorisations of $\ph$ and $\ph^n$, we find that $\ph_{n,\si} = c_n (\ph_{1,\si})^n$ for some constant $c_n$ with modulus $1$. It follows that $a_{\ph,n} = n a_{\ph,1}$ and $\sigma_{\ph,n} = n \sigma_{\ph,1}$. This, however, implies that $a_{\ph,1} = \tfrac{1}{n} a_{\ph,n} \ge - \tfrac{1}{n} a_g$ for all $n = 1, 2, \ldots$\,, and so $a_{\ph,1} \ge 0$. Similarly, the negative part of $\sigma_{\ph,1} = \tfrac{1}{n} \sigma_{\ph,n}$ is dominated by $\tfrac{1}{n} \sigma_g$ for any $n = 1, 2, \ldots\,$This is not possible if the negative part of $\sigma_{\ph,1}$ is non-zero, and therefore $\sigma_{\ph,1}$ is a~non-negative measure. We conclude that $\ph = \ph_{1,\bb}\,\ph_{1,\oo}\,\ph_{1,\si}$ is a bounded holomorphic function on $\CC_-$.

Since $\ph = f$ on $\CC_+$ and $f$ is a bounded holomorphic function on $\CC_+$, we have proved that $\ph$ is a bounded holomorphic function on $\CC \setminus A$. However, $A$ has zero Lebesgue measure (as a subset of $\R$). By Painlevé's theorem (see~\cite[Theorem~2.7]{bib:y15}),\linebreak$\ph$ extends to a bounded holomorphic function on $\CC$. This, in turn, implies that $\ph$ is constant, and so $\hat\mu_+$ is constant, contradicting the assumption that $\mu_+$ is a non-zero measure on $(0, \infty)$.
\end{proof}

\begin{proof}[Proof of Theorem~\ref{thm:wh}]
Suppose that $(\mu^{*n})_+ = (\nu^{*n})_+$ for $n = 1, 2,\,\ldots$ for some measures $\mu$ and $\nu$ such that $\mu_+$ and $\nu_+$ are non-zero measures. By Lemma~\ref{lem:mupm}, $\mu_+ = \nu_+$ and $(\mu_+^{*n} * \mu_-)_+ = (\nu_+^{*n} * \nu_-)_+$ for $n = 1, 2,\,\ldots\,$Let $\eta = \mu_+ + \mu_- - \nu_-$, so that $\eta_+ = \mu_+ = \mu_-$ and $\eta_- = \mu_- - \nu_-$. Then $(\eta_+^{*n} * \eta_-)_+ = 0$ for $n = 1, 2,\,\ldots\,$, and therefore, by Lemma~\ref{lem:wh}, either $\eta_+$ or $\eta_-$ is a~zero measure. Since $\eta_+ = \mu_+$ is a~non-zero measure, we must have $\eta_- = 0$, that is, $\mu_- = \nu_-$.
\end{proof}


\section{An error in~\cite{bib:u94}}

In~\cite{bib:u94} an analogue of Theorem~\ref{thm:wh} is given, with equality of $\mu^{*n}$ and $\nu^{*n}$ on $(-\infty, 0)$ rather than on $(0, \infty)$. In~\cite[Page~3001, line~16]{bib:u94}, it is claimed that the measures $\mu$ and $\nu$ satisfy~\cite[condition~(B) of Theorem~A]{bib:u94}, as a consequence of the results of~\cite[Section~11.2]{bib:lo77}. This reasoning would have been correct if the holomorphic extensions of $\hat{\mu}$ and $\hat{\nu}$ to the upper complex half-plane had been known to be continuous on the boundary. However, this is not verified in~\cite{bib:u94}.

More precisely, it is observed in~\cite{bib:u94} that $\hat{\mu} = (\hat{\chi}_2 - (\hat{\chi}_1)^2) / (2 \hat{\chi}_1)$ almost everywhere on~$\R$, where $\chi_1 = \mu - \nu$ and $\chi_2 = \mu^{*2} - \nu^{*2}$ are measures concentrated on $(0, \infty)$. Since $\hat{\chi}_1$ and $\hat{\chi}_2$ extend to holomorphic functions on $\CC_+$, $\hat{\mu}$ extends to a meromorphic function on $\CC_+$. Equality of $\mu^{*n}$ and $\nu^{*n}$ on $(-\infty, 0)$ for $n \ge 3$ is used only to show that the extension of $\hat{\mu}$ has no poles in $\CC_+$. However, the extension of $\hat{\mu}$ can have singularities near $\R$ and thus fail to satisfy~\cite[condition~(B) of Theorem~A]{bib:u94}.

To be specific, observe that $\hat{\mu}(z) = z^2 (z + i)^{-4} \exp(i/z)$ is the characteristic function of a measure $\mu$ on $\R$. Namely, $\mu$ is the convolution of $\tfrac{1}{6} x^3 e^{-x} \ind_{(0, \infty)}(x) dx$ and $\tfrac{1}{6} \, {_0F_1}(4; x) \ind_{(-\infty, 0)}(x) dx - \tfrac{1}{2} \delta_0(dx) - \delta_0'(dx) - \delta_0''(dx)$ (in the sense of distributions; ${_0F_1}$ is the hypergeometric function; we omit the details). Clearly, $\hat{\mu}$ extends holomorphically to the upper complex half-plane, but this extension is not continuous on the boundary, and thus $\mu$ does not satisfy~\cite[condition~(B) of Theorem~A]{bib:u94}. Furthermore, $\hat{\mu}(z)$ is the ratio of two characteristic functions of finite measures supported in $[0, \infty)$: $z^4 / (z + i)^8$ and $z^2 (z + i)^{-4} \exp(-i/z)$.

The author of the present article was not able to correct the error in~\cite{bib:u94}. The proof given above uses a related, but essentially different idea.


%\medskip

\subsection*{Acknowledgments}
I thank Lo\"ic Chaumont for numerous discussions on the subject of the article and encouragement. I thank Jan Rosiński for letting me know about reference~\cite{bib:ou90}.\linebreak I thank Alexander Ulanovski\u{\i} for discussions on references~\cite{bib:ou90,bib:u94}. The main part of this article was written during the \emph{$39^{\rm th}$ Conference on Stochastic Processes and their Applications} in Moscow.



\begin{thebibliography}{MMMMMM}
\bibitem[CD20]{bib:cd17} Chaumont, Loïc; Doney, Ron On distributions determined by their upward, space-time Wiener–Hopf factor, J. Theor. Probab., Volume 33 (2020) no. 2, pp. 1011-1033 \href{https://ahl.centre-mersenne.org/item/AHL_2020__3__1389_0/#bib:cd17}{Publisher reference}.
\bibitem[Con73]{bib:c73} Conway, John B. Functions of one complex variable, Graduate Texts in Mathematics, 11, Springer, 1973 \href{https://ahl.centre-mersenne.org/item/AHL_2020__3__1389_0/#bib:c73}{Publisher reference}.
\bibitem[Gar07]{bib:g07} Garnett, John Brady Bounded Analytic Functions, Graduate Texts in Mathematics, 236, Springer, 2007 \href{https://ahl.centre-mersenne.org/item/AHL_2020__3__1389_0/#bib:g07}{Publisher reference}.
\bibitem[LMS76]{bib:lms76} Letac, Gérard G.; Mazet, Pierre; Schiffman, Gérard A note on left-continuous random walks, J. Appl. Probab., Volume 13 (1976) no. 4, pp. 814-817 \href{https://ahl.centre-mersenne.org/item/AHL_2020__3__1389_0/#bib:lms76}{Publisher reference}.
\bibitem[LO77]{bib:lo77} Linnik, Yuriĭ V.; Ostrovskiĭ, Iosif V. Decomposition of random variables and vectors, Translations of Mathematical Monographs, 48, American Mathematical Society, 1977 \href{https://ahl.centre-mersenne.org/item/AHL_2020__3__1389_0/#bib:lo77}{Publisher reference}.
\bibitem[Mas09]{bib:m09} Mashreghi, Javad Representation theorems in Hardy spaces, London Mathematical Society Student Texts, Cambridge University Press, 2009 \href{https://ahl.centre-mersenne.org/item/AHL_2020__3__1389_0/#bib:m09}{Publisher reference}.
\bibitem[Ost85]{bib:o85} Ostrovskiĭ, Iosif V. Generalization of the Titchmarsh convolution theorem and the complex-valued measures uniquely determined by their restrictions to a half-line, Stability Problems for Stochastic Models (Uzhgorod, 1984) (Lecture Notes in Mathematics), Volume 1155, Springer, 1985, pp. 256-283 \href{https://ahl.centre-mersenne.org/item/AHL_2020__3__1389_0/#bib:o85}{Publisher reference}.
\bibitem[OU90]{bib:ou90} Ostrovskiĭ, Iosif V.; Ulanovskiĭ, Alexander M. Probability distributions and Borel measures uniquely determined by their restrictions to a half-space, Probability theory and mathematical statistics. Vol. II (Vilnius, 1989) (Grigelionos, B.; Prohorov, Yu V.; Sazonov, V. V.; Statulevičius, V., eds.), Mokslas, 1990, pp. 278-287 \href{https://ahl.centre-mersenne.org/item/AHL_2020__3__1389_0/#bib:ou90}{Publisher reference}.
\bibitem[Ula90]{bib:u90} Ulanovskiĭ, Alexander M. Unique determination of the convolutions of measures in R m , m\(\geqslant\)2, by their restriction to a set, J. Math. Sci., Volume 49 (1990) no. 6, pp. 1298-1301 \href{https://ahl.centre-mersenne.org/item/AHL_2020__3__1389_0/#bib:u90}{Publisher reference}.
\bibitem[Ula92]{bib:u94} Ulanovskiĭ, Alexander M. On the determination of a measure by the restriction of its n-fold convolutions to a massive set, Teor. Funkts. Funkts. Anal. Prilozh., Volume 57 (1992), pp. 102-109 English translation in: J. Math. Sci. 77(1) (1995): 2997–3002 \href{https://ahl.centre-mersenne.org/item/AHL_2020__3__1389_0/#bib:u94}{Publisher reference}.
\bibitem[Vig01]{bib:v02} Vigon, Vincent Simplifiez vos Lévy en titillant la factorisation de Wiener–Hopf, Ph. D. Thesis, INSA de Rouen, (France) (2001) \href{https://ahl.centre-mersenne.org/item/AHL_2020__3__1389_0/#bib:v02}{Publisher reference}.
\bibitem[You15]{bib:y15} Younsi, M. On removable sets for holomorphic functions, EMS Surv. Math. Sci., Volume 2 (2015) no. 2, pp. 219-254 \href{https://ahl.centre-mersenne.org/item/AHL_2020__3__1389_0/#bib:y15}{Publisher reference}.
\end{thebibliography}
\end{document}
